%! Change in Tandem — Unit 1, Lesson 1.2 — Guided Notes
\documentclass[10pt]{article}
\usepackage{precalculus-article}
\usepackage{precalculus-boxes}
\newcommand{\TallMath}[1]{$\displaystyle #1\rule[-1.4em]{0pt}{3.2em}$}

\begin{document}

\pageheader{Unit 1, Lesson 1.2}{Guided Notes}

\vspace{0.12in}

\begin{objectivebox}
By the end of this lesson, I will be able to\ldots
\begin{enumerate}[itemsep=2pt]
  \item Identify intervals where a function is \blank{2.8cm} or
        \blank{2.8cm}.
  \item Decide whether a graph is concave \blank{1.4cm} or concave
        \blank{1.4cm} by tracking the rate of change.
  \item Identify the \blank{1.8cm} of a function and explain what they mean
        in context.
\end{enumerate}
\end{objectivebox}

\vspace{0.1in}

\boxguard
\begin{vocabbox}
Fill in each term as we define it together.
\par\vspace{2pt}
\noindent\textbf{\textcolor{plum}{Increasing:}}\\[1pt]
\writeline\\[3pt]
\noindent\textbf{\textcolor{plum}{Decreasing:}}\\[1pt]
\writeline\\[3pt]
\noindent\textbf{\textcolor{plum}{Concave up:}}\\[1pt]
\writeline\\[3pt]
\noindent\textbf{\textcolor{plum}{Concave down:}}\\[1pt]
\writeline\\[3pt]
\noindent\textbf{\textcolor{plum}{Zero:}}\\[1pt]
\writeline
\end{vocabbox}

\vspace{0.1in}

\boxguard
\begin{hookbox}
A car pulls away from a stop sign. Two quantities start changing at once:
its \textit{speed} climbs quickly at first, then levels off near the limit;
its \textit{distance} from the sign grows slowly at first, then faster and
faster.

Both quantities are increasing. In words, what makes the two stories
different?
\writelines{2}
\end{hookbox}

\vspace{0.1in}

\boxguard
\begin{notesbox}{1. Two Quantities in Tandem}
A hose fills a tank. As time passes, the volume of water responds --- the two
quantities \textbf{vary in tandem}, linked by the function rule.

\smallskip
\renewcommand{\arraystretch}{1.4}
\begin{tabular}{|c|c|c|c|c|}
\hline
$t$ (min) & 0 & 2 & 4 & 6 \\
\hline
$V$ (gal) & 0 & 8 & 16 & 24 \\
\hline
\end{tabular}
\smallskip

\begin{itemize}[itemsep=4pt]
  \item The input (independent) quantity is \blank{2.6cm}; the output
        (dependent) quantity is \blank{2.6cm}.
  \item Every 2 minutes the tank gains \blank{1.4cm} gallons --- a perfectly
        steady climb.
  \item As $t$ increases, $V$ \blank{2.2cm} --- and a graph would show it
        rising left to right.
\end{itemize}
\end{notesbox}

\vspace{0.1in}

\boxguard[30]
\begin{notesbox}{2. Increasing and Decreasing}
$f$ is \textbf{increasing} on an interval when larger inputs give
\blank{2.2cm} outputs: whenever $a < b$, we get $f(a)$ \blank{1cm} $f(b)$.

$f$ is \textbf{decreasing} on an interval when larger inputs give
\blank{2.2cm} outputs: whenever $a < b$, we get $f(a)$ \blank{1cm} $f(b)$.

\smallskip
\textbf{From a graph.} The graph of $f$ is below.

\begin{center}
\begin{tikzpicture}[x=0.6cm, y=0.55cm]
  \draw[linegray, very thin, step=1] (0,0) grid (8,5);
  \draw[->, thick] (0,0) -- (8.6,0) node[below right] {\small $x$};
  \draw[->, thick] (0,0) -- (0,5.6) node[above] {\small $f(x)$};
  \foreach \x in {1,2,3,4,5,6,7,8} \node[below] at (\x,0) {\small \x};
  \foreach \y in {1,2,3,4,5} \node[left] at (0,\y) {\small \y};
  \draw[very thick, plum] (0,1) -- (3,5);
  \draw[very thick, plum] (3,5) -- (6,2);
  \draw[very thick, plum] (6,2) -- (8,4);
  \fill[plum] (0,1) circle (2.5pt);
  \fill[plum] (8,4) circle (2.5pt);
\end{tikzpicture}
\end{center}

\begin{itemize}[itemsep=4pt]
  \item Increasing for \blank{1cm} $< x <$ \blank{1cm} \ and for
        \blank{1cm} $< x <$ \blank{1cm}
  \item Decreasing for \blank{1cm} $< x <$ \blank{1cm}
  \item Intervals are always named by the \blank{1.6cm} ($x$-values) ---
        never by the outputs.
\end{itemize}

\textbf{Watch out:} decreasing does \textbf{not} mean negative. From $x = 3$
to $x = 6$ the outputs fall from 5 to 2, yet every output stays
\blank{1.8cm}.
\end{notesbox}

\vspace{0.1in}

\boxguard[30]
\begin{notesbox}{3. Concavity --- How the Change Changes}
Refill the tank, but leave a drain half-open so it fills fast at first:

\smallskip
\renewcommand{\arraystretch}{1.4}
\begin{tabular}{|c|c|c|c|c|}
\hline
$t$ (min) & 0 & 2 & 4 & 6 \\
\hline
$V$ (gal) & 0 & 14 & 20 & 24 \\
\hline
\end{tabular}
\smallskip

\begin{itemize}[itemsep=4pt]
  \item The gains are \blank{1cm} , \blank{1cm} , \blank{1cm} gallons ---
        still increasing, but the rate of change is \blank{2.2cm}.
  \item Rate of change decreasing $\Rightarrow$ \textbf{concave}
        \blank{1.4cm}.
\end{itemize}

A savings account grows the opposite way:

\smallskip
\renewcommand{\arraystretch}{1.4}
\begin{tabular}{|c|c|c|c|c|}
\hline
$m$ (months) & 0 & 1 & 2 & 3 \\
\hline
$S$ (\$) & 100 & 120 & 145 & 175 \\
\hline
\end{tabular}
\smallskip

\begin{itemize}[itemsep=4pt]
  \item The gains are \blank{1cm} , \blank{1cm} , \blank{1cm} dollars ---
        the rate of change is \blank{2.2cm}.
  \item Rate of change increasing $\Rightarrow$ \textbf{concave}
        \blank{1.4cm}.
\end{itemize}

\textbf{The two bends.} Label each curve \textit{concave up} or
\textit{concave down}.

\begin{center}
\begin{tikzpicture}[x=0.8cm, y=0.6cm]
  \draw[->, thick] (0,0) -- (3.4,0);
  \draw[->, thick] (0,0) -- (0,2.8);
  \draw[very thick, plum] (0,0.2) .. controls (1.8,0.5) .. (3,2.4);
  \node[below] at (1.6,-0.15) {\small slow, then fast};
\end{tikzpicture}
\hspace{0.6in}
\begin{tikzpicture}[x=0.8cm, y=0.6cm]
  \draw[->, thick] (0,0) -- (3.4,0);
  \draw[->, thick] (0,0) -- (0,2.8);
  \draw[very thick, plum] (0,0.2) .. controls (1.2,2.1) .. (3,2.4);
  \node[below] at (1.6,-0.15) {\small fast, then slow};
\end{tikzpicture}
\end{center}

Left curve: concave \blank{1.2cm} \qquad Right curve: concave \blank{1.2cm}

\smallskip
\textbf{Independent questions.} ``Increasing or decreasing?'' and ``concave
up or down?'' are separate calls --- the half-open-drain tank is
\blank{2.2cm} \textit{and} concave \blank{1.4cm} at the same time.
\end{notesbox}

\vspace{0.1in}

\boxguard[30]
\begin{notesbox}{4. Zeros}
A \textbf{zero} of a function is an input where the output is \blank{1cm};
the graph crosses or touches the horizontal axis there.

\smallskip
\textbf{From a graph.} The graph of $g$ is below.

\begin{center}
\begin{tikzpicture}[x=0.6cm, y=0.55cm]
  \draw[linegray, very thin, step=1] (0,0) grid (8,4);
  \draw[->, thick] (0,0) -- (8.6,0) node[below right] {\small $x$};
  \draw[->, thick] (0,0) -- (0,4.6) node[above] {\small $g(x)$};
  \foreach \x in {1,2,3,4,5,6,7,8} \node[below] at (\x,0) {\small \x};
  \foreach \y in {1,2,3,4} \node[left] at (0,\y) {\small \y};
  \draw[very thick, plum] (0,4) -- (2,0);
  \draw[very thick, plum] (2,0) -- (4,3);
  \draw[very thick, plum] (4,3) -- (6,0);
  \draw[very thick, plum] (6,0) -- (8,2);
  \fill[plum] (0,4) circle (2.5pt);
  \fill[plum] (8,2) circle (2.5pt);
\end{tikzpicture}
\end{center}

\begin{itemize}[itemsep=4pt]
  \item The zeros of $g$ are $x = $ \blank{1.2cm} and $x = $ \blank{1.2cm}.
  \item In context a zero is an \textit{event}. If $h(8) = 0$ for the drone
        flight from Lesson 1.1, the sentence is:

        \vspace{0.05in}
        \writeline
\end{itemize}
\end{notesbox}

\vspace{0.1in}

\boxguard
\begin{practicebox}
\begin{enumerate}[itemsep=8pt]
  \item A tank is draining. The table shows the water remaining.

        \smallskip
        \renewcommand{\arraystretch}{1.4}
        \begin{tabular}{|c|c|c|c|c|}
        \hline
        $t$ (min) & 0 & 2 & 4 & 6 \\
        \hline
        $w$ (gal) & 30 & 24 & 16 & 6 \\
        \hline
        \end{tabular}
        \smallskip

        \begin{enumerate}[label=(\alph*), itemsep=4pt]
          \item Is $w$ increasing or decreasing? \blank{2.6cm}
          \item The drops are 6, then \blank{1cm} , then \blank{1cm} gallons
                --- the water is draining faster and faster, so $w$ is
                concave \blank{1.4cm}.
          \item If $w(7) = 0$, what does that zero mean here?

                \vspace{0.05in}
                \writeline
        \end{enumerate}
\end{enumerate}
\end{practicebox}

\end{document}
